% % Detailed LaTeX document for Lab 1: Creating a Customer‑Support Agent Prototype
% % This standalone chapter provides an in‑depth exploration of the first
% % laboratory exercise in the Amazon Bedrock AgentCore workshop.  It is
% % written to be self‑contained, with explanations of every step and
% % extensive code listings.  Citations to the original workshop
% % materials are provided as footnotes; however, this document removes
% % any special characters that could prevent LaTeX compilation.

% \documentclass[12pt]{article}

% Packages for layout and formatting
% \usepackage{geometry}
% \geometry{margin=1in}
% \usepackage{setspace}
% \setstretch{1.2}
% \usepackage{amsmath, amssymb}
% \usepackage{graphicx}
% \usepackage{hyperref}
% \hypersetup{hidelinks}
% \usepackage{xcolor}
% \usepackage{listings}

% Define colours for code listings
\definecolor{keyword}{RGB}{0,0,128}
\definecolor{comment}{RGB}{0,96,0}
\definecolor{string}{RGB}{128,0,0}
\lstset{
  language=Python,
  basicstyle=\ttfamily\small,
  keywordstyle=\color{keyword}\bfseries,
  commentstyle=\color{comment}\itshape,
  stringstyle=\color{string},
  showstringspaces=false,
  frame=single,
  framerule=0.5pt,
  numbers=left,
  numberstyle=\tiny,
  numbersep=5pt
}

% \begin{document}

% \title{Lab 1: Creating a Customer‑Support Agent Prototype with Amazon Bedrock AgentCore}
% \author{}
% \date{}
% \maketitle

% \tableofcontents

\chapter{Creating a Customer‑Support Agent Prototype with Amazon Bedrock AgentCore}
\label{chap:Creating a Customer‑Support Agent Prototype with Amazon Bedrock AgentCore}
\section{Introduction and Goals}

This chapter guides you through building a basic customer‑support agent using
the Amazon Bedrock AgentCore and Strands Agents frameworks.  The purpose of
Lab 1 is to familiarise you with the fundamental building blocks of an AI
agent: creating specialised tools, crafting a system prompt, configuring a
foundation model and assembling these components into a functional agent.

The prototype agent you will build answers common customer questions about
consumer‑electronics products such as smartphones, laptops, headphones and
monitors.  It can tell a customer how long they have to return a product,
provide technical specifications and warranty information, search the web for
general questions and consult a curated knowledge base for troubleshooting
guidance.  At this stage the agent runs locally and has several
limitations—there is no persistent memory, no identity or authentication, and
no centralised tool sharing.  Later labs in the series will address these
concerns by adding memory, gateway and identity support, serverless
deployment, online evaluation and a customer‑facing frontend.

\section{Environment Setup}

Before writing any code, make sure your environment is properly configured.
Amazon Bedrock currently requires Python 3.10 or later, and the Strands and
AgentCore libraries rely on modern language features.  The following steps
outline a typical setup:

\begin{enumerate}
  \item \textbf{Install Python 3.10+ and create a virtual environment.}  On
        most systems you can create and activate a virtual environment with

        \begin{lstlisting}[language=bash]
python3 -m venv venv
source venv/bin/activate
        \end{lstlisting}

  \item \textbf{Install dependencies.}  Create a file named
        \texttt{requirements.txt} containing the following packages:

        \begin{lstlisting}[language=bash]
strands-agents>=0.7.0
bedrock-agentcore-starter-toolkit>=0.3.0
duckduckgo-search>=4.0.0
boto3>=1.29.0
        \end{lstlisting}

        Then install them using pip:

        \begin{lstlisting}[language=bash]
pip install -r requirements.txt
        \end{lstlisting}

  \item \textbf{Configure your AWS credentials.}  Run
        \texttt{aws configure} and provide your Access Key, Secret Key and
        region.  The agent will interact with Amazon Bedrock, Systems
        Manager (for Parameter Store) and S3, so valid credentials are
        essential.  You should also ensure that you have access to an
        Anthropic Claude model on Bedrock, such as \texttt{anthropic.claude-v2}.  If you
        do not yet have access, request it through the AWS console.

  \item \textbf{Clone the workshop repository.}  The official code is
        available on GitHub.  Clone it and change into the Lab 1
        directory so you can follow along with the notebook and raw files:

        \begin{lstlisting}[language=bash]
git clone https://github.com/awslabs/amazon-bedrock-agentcore-samples.git
cd amazon-bedrock-agentcore-samples/01-tutorials/09-AgentCore-E2E
        \end{lstlisting}

\end{enumerate}

With your environment prepared, you are ready to implement the tools that
power the agent.

\section{Implementing Custom Tools}

In Strands Agents, a \emph{tool} is a callable function decorated with
\texttt{@tool}.  Each tool must accept JSON‑serialisable inputs and return
JSON‑serialisable outputs.  The decorator registers the function with the
agent so the language model can invoke it when appropriate.  It is best
practice to provide a clear docstring that explains what the tool does and
when it should be used.

Lab 1 defines four tools:

\begin{itemize}
  \item \textbf{Return‑policy tool} — returns the return and warranty policy for
        a given product category.
  \item \textbf{Product information tool} — provides technical
        specifications, features, compatibility and support contact for a
        specific product type.
  \item \textbf{Web search tool} — performs an internet search using
        DuckDuckGo when the agent cannot answer from its own knowledge.
  \item \textbf{Technical support tool} — retrieves troubleshooting steps from
        a knowledge base stored in Amazon Bedrock.
\end{itemize}

Each tool is described in detail below, including the full Python
implementation and a discussion of design choices.

\subsection{Return‑Policy Tool}

Customers often ask about return deadlines, acceptable product condition and
warranty coverage.  The \texttt{get\_return\_policy} function encapsulates
these rules.  It uses a dictionary to map product categories to structured
policy information and returns a formatted, human‑readable string.  Unknown
categories trigger a polite error message.

\begin{lstlisting}[language=Python]
from strands import tool

@tool
def get_return_policy(category: str) -> str:
    """Return policy for the given product category.

    Parameters
    ----------
    category : str
        The type of product (e.g., 'smartphones', 'laptops', 'accessories').

    Returns
    -------
    str
        A multi-line description of return windows, conditions, process,
        refund timing, shipping and warranty.  If the category is unknown
        the function returns a helpful message.
    """
    # Define return policies for each category
    policies = {
        "smartphones": {
            "window": 30,
            "condition": "Unopened or like new",
            "process": "Return via store or mail",
            "refund_time": "5 business days",
            "shipping": "Free return shipping",
            "warranty": "1-year limited",
        },
        "laptops": {
            "window": 14,
            "condition": "Factory sealed or like new",
            "process": "Return to store or mail",
            "refund_time": "10 business days",
            "shipping": "Customer pays return shipping",
            "warranty": "2-year extended",
        },
        "accessories": {
            "window": 90,
            "condition": "Unused and unopened",
            "process": "Return to store",
            "refund_time": "3 business days",
            "shipping": "Free return shipping",
            "warranty": "1-year limited",
        },
    }

    # Look up the policy for the requested category
    policy = policies.get(category.lower())
    if policy is None:
        return (f"We do not have a return policy for the category '{category}'."
                " Please check the product documentation or contact support.")

    # Build a formatted response from the policy fields
    lines = [
        f"Return window: {policy['window']} days",
        f"Condition: {policy['condition']}",
        f"Process: {policy['process']}",
        f"Refund time: {policy['refund_time']}",
        f"Shipping: {policy['shipping']}",
        f"Warranty: {policy['warranty']}",
    ]
    return "\n".join(lines)
\end{lstlisting}

\paragraph{Discussion.}  Note how the tool adheres to the single
responsibility principle: it performs exactly one operation—fetching and
formatting return policy information.  The dictionary structure makes it
trivial to add new categories or edit existing policies without changing the
control flow.  The docstring provides a clear specification so that both the
developer and the language model understand how to use the tool.  Finally,
the function normalises the category to lowercase to handle inputs such as
``SmartPhones'' or ``LAPTOPS'' gracefully.

\subsection{Product‑Information Tool}

The next tool, \texttt{get\_product\_info}, provides technical and support
information about various product types.  Like the return‑policy tool, it
organises data into nested dictionaries.  Each product entry contains a
warranty period, a specification sub‑dictionary, a list of distinctive
features, a compatibility field and a support email address.  The function
constructs a consolidated summary from these fields.

\begin{lstlisting}[language=Python]
@tool
def get_product_info(product: str) -> str:
    """Provide technical specifications, features and support information.

    Parameters
    ----------
    product : str
        The type of product (e.g., 'smartphones', 'laptops', 'headphones',
        'monitors').

    Returns
    -------
    str
        A multi-line description containing warranty, specifications,
        features, compatibility and support contact details.  Unknown
        products return a friendly message.
    """
    products = {
        "smartphones": {
            "warranty": "1 year",
            "specs": {
                "cpu": "Octa-core 3.2 GHz",
                "ram": "8 GB",
                "storage": "128 GB",
                "display": "6.4\" AMOLED",
            },
            "features": ["fast charging", "water resistant", "5G support"],
            "compatibility": "USB‑C",
            "support": "support@example.com",
        },
        "laptops": {
            "warranty": "2 years",
            "specs": {
                "cpu": "Intel i7",
                "ram": "16 GB",
                "storage": "512 GB SSD",
                "display": "15\" IPS",
            },
            "features": ["backlit keyboard", "fingerprint reader",
                          "Thunderbolt 4"],
            "compatibility": "USB‑C",
            "support": "support@example.com",
        },
        "headphones": {
            "warranty": "1 year",
            "specs": {
                "drivers": "40 mm",
                "connectivity": "Bluetooth 5.2",
                "battery": "30 hours",
            },
            "features": ["noise cancellation", "wireless", "fast pairing"],
            "compatibility": "Bluetooth devices",
            "support": "support@example.com",
        },
        "monitors": {
            "warranty": "2 years",
            "specs": {
                "size": "27\"",
                "resolution": "2560×1440",
                "panel": "IPS",
                "refresh_rate": "144 Hz",
            },
            "features": ["HDR", "USB hub", "adjustable stand"],
            "compatibility": "HDMI/DisplayPort",
            "support": "support@example.com",
        },
    }

    info = products.get(product.lower())
    if info is None:
        return (f"Sorry, I don't recognise the product '{product}'."
                " Please ask about smartphones, laptops, headphones or monitors.")

    # Build a summary from the nested dictionary
    spec_details = ', '.join(f"{k}: {v}" for k, v in info['specs'].items())
    summary_lines = [
        f"Warranty: {info['warranty']}",
        f"Specs: {spec_details}",
        f"Features: {', '.join(info['features'])}",
        f"Compatibility: {info['compatibility']}",
        f"Support: {info['support']}",
    ]
    return "\n".join(summary_lines)
\end{lstlisting}

\paragraph{Discussion.}  The product‑information tool showcases how to
organise complex data in dictionaries and lists.  Grouping specification
fields under a \texttt{specs} sub‑dictionary keeps the top‑level structure
clean and makes it easy to iterate over the key–value pairs.  The function
creates a concise, readable summary by joining the lists and dictionaries
into strings.  If you wish to add more products, simply insert another
entry in the \texttt{products} dictionary.

\subsection{Web‑Search Tool}

Sometimes customers ask questions that fall outside the curated product data
and knowledge base.  For example, ``When is the next version of our
flagship phone expected to launch?''  The agent can use a web search to
provide a general answer.  The \texttt{web\_search} tool leverages the
\texttt{duckduckgo\_search} library to perform an internet search and
returns a list of dictionaries containing the title, URL and snippet for
each result.  If network errors or rate limits occur, the tool returns a
single entry describing the issue.

\begin{lstlisting}[language=Python]
from duckduckgo_search import DDGS

@tool
def web_search(query: str) -> list[dict]:
    """Search the web using DuckDuckGo.

    Parameters
    ----------
    query : str
        The search query provided by the customer.

    Returns
    -------
    list of dict
        A list of search results.  Each entry contains keys 'title',
        'url' and 'snippet'.  If an error occurs, a single dictionary
        with a message is returned.
    """
    try:
        ddgs = DDGS()
        results = ddgs.search(query, max_results=5)
        return [
            {"title": r.title, "url": r.href, "snippet": r.body}
            for r in results
        ]
    except Exception:
        # Rate limited or network error
        return [
            {
                "title": "Search error",
                "url": "",
                "snippet": (
                    "The web search service is currently unavailable. "
                    "Please try again later."
                ),
            }
        ]
\end{lstlisting}

\paragraph{Discussion.}  The return type of this tool is declared as
\texttt{list[dict]}, which helps both the type checker and the language
model.  Each dictionary field is explicit, ensuring that the agent can
render the search results clearly.  The try/except block catches any
exceptions raised by the \texttt{DDGS} API (e.g., network errors or rate
limits) and returns an informative fallback result.

\subsection{Technical‑Support Tool}

The final tool in Lab 1 fetches troubleshooting content from a knowledge
base.  The knowledge base is created and populated separately (see
Section~\ref{sec:kb}); the tool simply retrieves the top relevant
document for a customer query.  It reads the knowledge base and data
source identifiers from AWS Systems Manager Parameter Store and invokes
the Bedrock knowledge base retrieval API via Strands.

\begin{lstlisting}[language=Python]
import boto3
from strands import retrieve

@tool
def get_technical_support(query: str) -> str:
    """Retrieve troubleshooting information from the knowledge base.

    Parameters
    ----------
    query : str
        The customer's support question.

    Returns
    -------
    str
        The text of the most relevant document in the knowledge base, or a
        default message if nothing is found.
    """
    ssm = boto3.client("ssm")
    # Look up identifiers stored by the knowledge base ingestion script
    kb_id = ssm.get_parameter(Name="/customer_support/kb_id")["Parameter"]["Value"]
    ds_id = ssm.get_parameter(Name="/customer_support/ds_id")["Parameter"]["Value"]
    # Query the knowledge base using the Strands retrieve function
    results = retrieve(
        knowledge_base_id=kb_id,
        input=query,
        retrieval_params={"top_k": 3, "format": "text"},
    )
    # Return the first result's text if available
    return results.get("text", "No relevant answer found in the knowledge base.")
\end{lstlisting}

\paragraph{Discussion.}  This tool demonstrates how to integrate with
external services.  The SSM client retrieves configuration values that
identify the knowledge base and its data source.  The Strands
\texttt{retrieve} function performs a vector search against the knowledge
base and returns a dictionary.  Adjust the \texttt{top\_k} parameter to
return more or fewer documents, and experiment with the \texttt{format}
field to obtain plain text or citations.

\section{Preparing the Knowledge Base}
\label{sec:kb}

To answer hardware troubleshooting questions, the agent must have access to
relevant documentation.  In the workshop, a small set of PDF manuals and
troubleshooting guides are stored in an Amazon S3 bucket.  Before you can
retrieve documents with the technical‑support tool, you need to download
these files and ingest them into a Bedrock knowledge base.

The ingestion process is handled by a helper script in the workshop
repository.  Below is a simplified version of that script.  It first
downloads all objects under a given prefix from an S3 bucket into a local
folder and then triggers a Bedrock ingestion job.  Finally, it polls the
status of the job until it finishes.

\begin{lstlisting}[language=Python]
import boto3
import os
import time

# Configure your S3 bucket and prefix
bucket_name = "your-knowledge-base-bucket"
prefix = "docs/"

# Create clients
s3 = boto3.client("s3")
agent = boto3.client("bedrock-agent")

# Create the output directory
target_dir = "knowledge_base_data/"
os.makedirs(target_dir, exist_ok=True)

# Download each object in the prefix
objects = s3.list_objects_v2(Bucket=bucket_name, Prefix=prefix).get("Contents", [])
for obj in objects:
    key = obj["Key"]
    filename = key.split("/")[-1]
    s3.download_file(bucket_name, key, os.path.join(target_dir, filename))

# Start an ingestion job
response = agent.start_ingestion_job(
    knowledgeBaseId=kb_id,
    dataSourceId=ds_id,
)
job_id = response["ingestionJob"]["jobId"]
print(f"Started ingestion job {job_id}")

# Poll the job status until complete
while True:
    status = agent.get_ingestion_job(
        knowledgeBaseId=kb_id,
        dataSourceId=ds_id,
        jobId=job_id,
    )["ingestionJob"]["status"]
    if status in ("COMPLETED", "FAILED"):
        break
    print(f"Job {job_id} status = {status}; waiting…")
    time.sleep(10)

print(f"Ingestion job completed with status {status}")
\end{lstlisting}

After the ingestion job finishes successfully, the knowledge base ID
\texttt{kb\_id} and data source ID \texttt{ds\_id} should be stored in
Parameter Store.  The technical‑support tool reads these parameters at
runtime, ensuring that the agent always queries the correct knowledge base.

\section{Designing the System Prompt}

A system prompt defines the agent’s personality and rules of engagement.
Large language models rely heavily on this prompt to decide whether to call a
tool, answer directly, or ask for clarification.  A good system prompt is
explicit about the tools available and instructs the model to use them when
necessary rather than fabricating answers.

Here is an example prompt used in the workshop:

\begin{quote}
\textbf{You are a professional customer support agent for an electronics
store.}  You have the following tools: \texttt{get\_return\_policy},
\texttt{get\_product\_info}, \texttt{web\_search} and
\texttt{get\_technical\_support}.  Use these tools to answer customer
questions.  When you do not know the answer, call the appropriate tool
instead of guessing.  Be polite, concise and professional.
\end{quote}

Feel free to customise the wording to reflect your brand’s tone of voice.  A
well‑crafted prompt can dramatically improve tool usage and reduce the risk
of hallucinated answers.

\section{Assembling the Agent}

With the tools and prompt defined, you can construct the agent object.  The
\texttt{strands.Agent} class requires a system prompt, a list of tools and
 a model configuration.  The Bedrock model is configured via a small helper
class that specifies the model identifier, temperature and region.  In this
example we use Anthropic Claude v2 with deterministic settings.

\begin{lstlisting}[language=Python]
from strands import Agent
from bedrock_agent_toolkit.models.bedrock import BedrockModel

# Define the system prompt as a global constant
SYSTEM_PROMPT = (
    "You are a professional customer support agent for an electronics store.\n"
    "You have the following tools: get_return_policy, get_product_info, web_search, get_technical_support.\n"
    "Use these tools to answer customer questions. When you don't know the answer, call the appropriate tool instead of guessing.\n"
    "Be polite, concise and professional."
)

# Configure the underlying foundation model
model = BedrockModel(
    model_id="anthropic.claude-v2",
    temperature=0.0,
    region="us-east-1",
)

# Instantiate the agent with the prompt, tools and model
agent = Agent(
    system_prompt=SYSTEM_PROMPT,
    tools=[get_return_policy, get_product_info, web_search, get_technical_support],
    model=model,
)

# Invoke the agent with a customer question
response = agent.invoke("What is the warranty period for laptops?")
print(response)
\end{lstlisting}

\paragraph{What happens during invocation?}  The \texttt{agent.invoke}
method performs several steps under the hood:

\begin{enumerate}
  \item \emph{Prompting the model.}  The agent sends the system prompt
        along with the customer's question to the foundation model.
  \item \emph{Tool selection.}  The model analyses the query and decides whether
        to call a tool.  If the query requires data not contained in the
        model’s internal knowledge, it will emit a tool call with the
        appropriate arguments.
  \item \emph{Tool execution.}  The agent detects the tool call,
        executes the corresponding Python function and captures the result.
  \item \emph{Response composition.}  The model receives the tool output and
        integrates it into a final answer.  This step may be repeated if
        multiple tools are needed.
  \item \emph{Delivery.}  The agent returns the generated response to the
        caller.  In a local notebook this simply prints the answer; in
        production it could be sent to a web client via an API.
\end{enumerate}

During development you can inspect the intermediate events and tool calls
using logging or the built‑in Strands tracing hooks.  Subsequent labs
introduce robust observability features via AgentCore runtime.

\section{Local Testing}

Testing your agent in an interactive environment helps uncover issues before
deploying.  Here are some example interactions you can try in a Python
shell or notebook:

\begin{itemize}
  \item \textbf{Return policy:}
        \begin{lstlisting}[language=Python]
agent.invoke("What's the return policy for smartphones?")
        \end{lstlisting}
        The agent should call \texttt{get\_return\_policy} and return a
        multi‑line description of the smartphone return policy.

  \item \textbf{Product information:}
        \begin{lstlisting}[language=Python]
agent.invoke("Tell me about your laptop models.")
        \end{lstlisting}
        This triggers \texttt{get\_product\_info} and responds with
        warranty, specs, features and compatibility details.

  \item \textbf{Technical support:}
        \begin{lstlisting}[language=Python]
agent.invoke("My laptop will not turn on.  What should I do?")
        \end{lstlisting}
        Here the agent retrieves the most relevant troubleshooting steps
        from the knowledge base via \texttt{get\_technical\_support}.

  \item \textbf{General enquiry:}
        \begin{lstlisting}[language=Python]
agent.invoke("When is the next version of the smartphone launching?")
        \end{lstlisting}
        Because this question is not covered by the internal data or
        knowledge base, the agent will call \texttt{web\_search} and
        summarise the search results.
\end{itemize}

While testing, pay attention to how the agent decides to call tools.  You
can alter the system prompt or tool docstrings to influence its behaviour.
For example, if the agent starts answering from memory instead of calling
the knowledge base, emphasise in the prompt that it must consult the tools.

\section{Limitations and Future Labs}

The prototype agent built in Lab 1 lays the groundwork for more advanced
capabilities, but it still has notable limitations:

\begin{itemize}
  \item \textbf{No memory.}  The agent cannot remember previous
        interactions, so each invocation is stateless.  Lab 2 introduces
        AgentCore memory to persist user preferences and conversation
        history.
  \item \textbf{No authentication or identity.}  Anyone can run the agent,
        and there is no way to limit tool access based on user roles.
        Lab 3 explains how to integrate Amazon Cognito and the
        AgentCore gateway to control access and share tools across agents.
  \item \textbf{Limited observability.}  Local testing offers little
        visibility into tool usage and model reasoning.  Lab 4 shows how to
        deploy the agent to the AgentCore runtime with built‑in tracing
        and monitoring.
  \item \textbf{Manual deployment and evaluation.}  Currently the agent
        runs only in your notebook.  Lab 5 introduces an online
        evaluation framework, and Lab 6 builds a customer‑facing web
        application.
\end{itemize}

Addressing these limitations will turn the simple prototype into a robust
customer‑support platform capable of handling real‑world workloads.


